% ===================================================================
%  नकसा नं. ८ — लवनी।
%
%  Traduction, faite en 2026, en bhojpouri d'aujourd'hui, sur l'ido
%  de texto/io. Regles de la colonne : voir 10-nakasa-01.tex. DEUX
%  scenes, deux numerotations qui repartent de 1 (t08-c1-*,
%  t08-c2-*). 40 blocs, 109 renvois, vingt-deux paires a retourner.
%  UNE note d'auteur, posee en tete.
%
%  ECART DECLARE : t08-c2-03-1. L'ido y compose le renvoi du gobelet
%  « \textsuperscript{13)} », sans la parenthese ouvrante ; le motif
%  ne se reconnait donc pas et renvoji.py laisse le bloc passer SANS
%  RIEN VERIFIER. La traduction retablit « (13) » comme les
%  quarante-six autres colonnes, et l'ordre a ete compte a la main :
%  3, 12, 14, 13, 15, 16.
%
%  TROIS OUTILS DE MOISSON, TROIS REPONSES, ET CHACUNE DIT COMMENT ON
%  MOISSONNE ICI. Le tableau 7 avait trouve le trou : pas de mot pour
%  la faux. Celui-ci complete la mesure sur les deux autres outils :
%
%    la faucille (25)   हँसुआ    la lame courte qu'on tient accroupi
%    la faucille (12)   दराँती   la lame dentee, l'autre geste
%    la faux            AUCUN    on ecrit « la grande faucille »
%    le fleau    (40)   AUCUN    on ne bat pas la gerbe : on la fait
%                                fouler par les boeufs — दाँवन
%
%  Le bhojpouri a DEUX mots la ou le francais n'en a qu'un, et
%  l'ido lui-meme distingue « falchilo » et « siklo ». Il n'en a
%  AUCUN pour la faux, ni pour le fleau : on ne bat pas la gerbe au
%  baton dans la plaine, on la fait fouler par les boeufs — दाँवन —,
%  et l'outil n'a donc jamais eu besoin d'exister. L'OUTIL QU'ON SE
%  SERT TOUS LES JOURS A DEUX NOMS ; CELUI QU'ON N'A JAMAIS EU N'EN
%  A PAS.
%
%  ET LE SEIGLE (43) PORTE LE NOM DE LA MOUTARDE, EXACTEMENT COMME EN
%  GUJARATI. La colonne gujaratie avait releve que « રાઈ » designe la
%  moutarde et non le seigle, et que la cereale a donc « un nom de
%  trop ». Le bhojpouri ecrit राई et le mot vaut moutarde ici aussi.
%  La raison est la meme des deux cotes de l'Inde : la graine
%  oleagineuse se seme partout, la cereale ne pousse nulle part, et
%  c'est le nom disponible qui a servi. DEUX LANGUES, DEUX
%  ECRITURES, LE MEME MOT ET LA MEME CONFUSION — ce qui la fait
%  passer du rang de curiosite gujaratie a celui de fait indien.
%
%  LE COQUELICOT (2) EST « जंगली पोस्ता », ET IL FALLAIT LE MARQUER.
%  पोस्ता tout court est le pavot a opium, et ce n'est pas une fleur
%  des champs dans cette region : Ghazipur, en plein pays bhojpouri,
%  abrite depuis 1820 la plus grande fabrique d'opium legale du
%  monde, et la culture y est sous licence d'Etat. Le petit pavot
%  rouge que Rochelle met dans son ble n'a rien a voir avec cela.
%  D'ou « जंगली पोस्ता », le pavot sauvage. La colonne gujaratie
%  avait REFUSE le coquelicot, faute de mot ; celle-ci en a un, et
%  doit le desarmer.
% ===================================================================

%%K t08-noto-1 noto
\VUnotes{143.2mm}{%
\textit{(*) काहेकि ई सबद साफ नइखे : ई सन के लवनी ह, अंगूर के लवनी ह,
भा सेब के लवनी ? सायद «} \VUgras{mesono} \textit{» लिखल नीक होखित, जवन
Mondo XIV, प. २४२ में सुझावल गइल। बाकिर अबहीं ले ई नया धातु अकादमी
मंजूर नइखे कइले आ इडो के चलन के हिसाब से बहस के अगोरत बा।}%
}

%%K t08-apar-1 apar
\VUsaut{5.0mm}
\VUtitre{15.0pt}{60}{नकसा नं. ८}
\VUsaut{3.0mm}
\VUcentre{13.2pt}{90}{\textit{पहिलका दृस्य।}}
\VUsaut{3.0mm}
\VUpk{10.2pt}{40}{लवनी (*)।}
\VUsaut{4.0mm}
\VUpk{10.2pt}{40}{देहात के नजारा।}
\VUsaut{3.0mm}

%%K t08-c1-01-1 p
१. --- \VUgras{गरमी} के साथे देहात के रूप फेरु बदल गइल। नार में डालल
बीया जम गइल ; हरियर पौधा माटी से निकलल आ ओहमें बाली आइल। दाना बनल,
फेरु पाकल, आ अबहीं खाली लवनी करे के बाकी बा।

%%K t08-c1-02-1 p
२. --- \VUgras{देहात के फूल} खिल गइल। \VUgras{नील फूल}
\textsuperscript{(1)}, \VUgras{जंगली पोस्ता} \textsuperscript{(2)} आ
\VUgras{तिलपुष्पी} \textsuperscript{(3)} गोहूँ के खेत के एके रंग में आपन
ताजा \VUgras{पंखुड़ी} के खुसनुमा भेद मिला देवेला।

%%K t08-c1-03-1 p
३. --- फर के गाछ पत्ता से भर गइल ; फर पाके लागल। छोट \VUgras{चेरी के
गाछ} \textsuperscript{(4)} सुंदर \VUgras{चेरी} \textsuperscript{(5)} से लदल
बा, जेकरा के लइकन बड़का चाव से तोड़ेला आ खाला।

%%K t08-c1-04-1 p
४. --- अबहीं जनावर पूरा दिन खुला हवा में बिता सकेला।

%%K t08-c1-04-2 p
\VUgras{मेमना} \textsuperscript{(6)} बड़ हो गइल आ सुंदर \VUgras{भेड़}
\textsuperscript{(7)} आ सुंदर \VUgras{मेढ़ा} \textsuperscript{(8)} बन गइल,
घन ऊन के साथे। ऊ लोग सांति से \VUgras{चरागाह} \textsuperscript{(9)} के
\VUgras{घास} चरत बा, जवन \VUgras{गुलबहार} से रंगल बा।

%%K t08-c1-04-3 p
जल्दिए \VUgras{एह} जनावरन के काटल जाई ताकि ओकर \VUgras{ऊन} मिले, जवना
से हमनी के कपड़ा के ऊनी थान बनेला।

%%K t08-tit-2 sub
\VUsaut{5.0mm}
\VUpk{10.2pt}{40}{गड़ेरिया।}
\VUsaut{3.4mm}

%%K t08-c1-05-1 p
५. --- \VUgras{गड़ेरिया} \textsuperscript{(10)} ओह लोग पर जादे धेयान नइखे
देत लागत। ओकरा खाली ओह आँधी के फिकिर बा जवन छितिज पर से गुजरत बा, आ
\VUgras{भेड़ के रखवार} \textsuperscript{(13)}, जे आपन \VUgras{तुमड़ी}
\textsuperscript{(14)} ओठ लगे ले जात बा, त आउरो बेफिकिर लागत बा।

%%K t08-c1-06-1 p
६. --- गरमी बहुते सतावत बा ; काहेकि \VUgras{कुकुर} \textsuperscript{(15)}
भी आपन पियास बुझावे आवत बा ; ऊ \VUgras{नाला} \textsuperscript{(16)} के
ठंढा पानी चाटत बा आ एगो बेचारा \VUgras{मेढ़क} \textsuperscript{(17)} के
डरा देत बा जे किनारे के घास में छिपल रहे। ई मेढ़क साफ पानी में कूद जात
बा जहाँ \VUgras{कमल} \textsuperscript{(18)} खिलल बा।

%%K t08-c1-07-1 p
७. --- \VUgras{खंजन} \textsuperscript{(19)} ओह \VUgras{सोता}
\textsuperscript{(20)} लगे नहात बा जवन दू गो चट्टान के बीच से फूटत बा, आ
\VUgras{कँटीला झाड़} \textsuperscript{(21)} आ \VUgras{सफेद काँट}
\textsuperscript{(22)} के नीचे छिप जात बा।

%%K t08-tit-3 sub
\VUsaut{5.0mm}
\VUpk{10.2pt}{40}{लवनी।}
\VUsaut{3.4mm}

%%K t08-c1-08-1 p
८. --- देहात के लइकन खातिर गोहूँ के लवनी बहुते मजेदार बेरा ह। ऊ लोग
खुसी से \VUgras{लोटत} \textsuperscript{(24)} बा \VUgras{पयार के ढेर}
\textsuperscript{(23)} पर, ऊ लोग \VUgras{लवनी करे वाला} \textsuperscript{(26)}
के मदद करत बा, आ जब ऊ लोग \VUgras{गोहूँ के खेत} \textsuperscript{(27)}
आपन \VUgras{हँसुआ} \textsuperscript{(25)} से काटत बा, जवन \VUgras{सान के
पाथर} \textsuperscript{(30)} पर नीक से तेज कइल बा, त लइकन गोहूँ बटोरत बा
आ ओकरा के \VUgras{पूला} \textsuperscript{(31)} भा \VUgras{मुट्ठा}
\textsuperscript{(32)} में बान्हत बा।

%%K t08-c1-09-1 p
९. --- कबो-कबो ओहमें से सबसे बड़का के इजाजत मिलेला कि ऊ
\VUgras{दराँती} \textsuperscript{(12)} से ओह \VUgras{सोना नियर डंठल}
\textsuperscript{(28)} के काटे जे दाना से भरल भारी \VUgras{बाली}
\textsuperscript{(29)} थामले बा ; आ छोटका लोग, \VUgras{रेवड़}
\textsuperscript{(33)} पर नजर राखत जे \VUgras{चरागाह}
\textsuperscript{(34)} में बड़का \VUgras{लिंडेन के गाछ}
\textsuperscript{(35)} के छाँह में चरत बा, आपन \VUgras{जाल} \textsuperscript{(36)} से सुंदर \VUgras{तितली}
पकड़त बा।

%%K t08-c1-09-2 p
कुछ लोग त \VUgras{गाड़ी} \textsuperscript{(37)} पर चढ़ के ओह पूला के
जमावत बा जवन ओह लोग के \VUgras{कँटिया} \textsuperscript{(38)} से थमावल
जात बा।

%%K t08-c1-10-1 p
१०. --- पूरा \VUgras{मैदान} \textsuperscript{(39)} पर, जहाँ
\VUgras{चंडूल} \textsuperscript{(41)} उड़े लागत बा, हलचल के राज बा। सब लोग
लवनी में लागल बा। बाकिर जब तक गरीब लोग पुरनका औजार चलावत बा —
\VUgras{हँसुआ, दराँती} आ \VUgras{दँवरी के डंडा} \textsuperscript{(40)} —
धनी जमींदार आपन गोहूँ भा आपन \VUgras{राई} \textsuperscript{(43)}
\VUgras{लवनी के मसीन} \textsuperscript{(42)} से कटवावेलें। ओकरा बाद, पूला
के खरिहान तक ढोवे के जरूरत ना पड़ेला, आ बड़का \VUgras{पूला के ढेर}
\textsuperscript{(44)} लगा देहल जाला।

%%K t08-c1-11-1 p
११. --- तब \VUgras{दँवरी के मसीन} \textsuperscript{(47)} ले आवल जाला,
जेकरा के \VUgras{चलत इंजन} \textsuperscript{(45)} \VUgras{पट्टा}
\textsuperscript{(46)} से चलावेला, आ एह तरह दाना \VUgras{पयार} से अलग हो
जाला, बिना ओह खेत के छोड़ले जहाँ ऊ बोवल गइल रहे।

%%K t08-tit-4 sub
\VUsaut{5.0mm}
\VUpk{10.2pt}{40}{आँधी।}
\VUsaut{3.4mm}

%%K t08-c1-12-1 p
१२. --- लवनी के बेरा अकसर \VUgras{आँधी-पानी} \textsuperscript{(53)}
आवेला। पहिले दूर से \VUgras{गरज} के आवाज सुनाला ; \VUgras{पछुआ हवा}
\textsuperscript{(54)} चले लागेला आ हाँफत \VUgras{जंगल} \textsuperscript{(49)}
के सबसे मोट गाछ के झुका देला। करिया \VUgras{बादर} \textsuperscript{(56)}
आसमान पर चढ़ आवेला ; ओकरा के \VUgras{बिजली} \textsuperscript{(55)} के
चौंधिआवत टेढ़ लकीर चीर देला। \VUgras{बरखा} आ कबो-कबो \VUgras{ओला} गिरे
लागेला, आ कटे खातिर तइयार फसल के दबा देला।

%%K t08-c1-13-1 p
१३. --- अकसर बिजली कवनो इमारत में आग लगा देला। जइसे पिछला साल
\VUgras{पवनचक्की} \textsuperscript{(50)} के छत \VUgras{राख}
\textsuperscript{(51)} हो गइल आ ओकर दू गो \VUgras{पाँख} टूट गइल।

%%K t08-c1-14-1 p
१४. --- कुछ घंटा के बाद फेरु सांति लवट आवेला। छितिज पर चमकत
\VUgras{इंद्रधनुष} \textsuperscript{(52)} लउकेला आ प्रकृति आपन ऊ जिनगी
फेरु सुरू करेले जवन कुछ पल खातिर रुक गइल रहे। जे देहाती लोग
\VUgras{मड़ई} \textsuperscript{(48)} में सरन लेले रहे, ऊ फेरु काम पर लाग
जाला आ हिम्मत से ओह नुकसान के भरपाई करे के कोसिस करेला जवन अभिए भइल बा।

%%K t08-tit-5 sub
\VUsaut{6.0mm}
\VUcentre{13.2pt}{90}{\textit{दुसरका दृस्य।}}
\VUsaut{3.6mm}

%%K t08-tit-6 sub
\VUpk{10.2pt}{40}{सिकार। --- अंगूर के लवनी।}
\VUsaut{1.4mm}
\VUpk{10.2pt}{40}{मछरी मारल।}
\VUsaut{4.0mm}

%%K t08-c2-01-1 p
१. --- \VUgras{अगस्त} के आखिर भा \VUgras{सितंबर} के बीच में, इलाका के
हिसाब से, सिकार के मौसम खुलेला। सूरज के पहिलका किरन से
\VUgras{छिपकली} \textsuperscript{(2)} आपन बिल से निकलबे ना कइले रहेली कि
\VUgras{सिकारी} \textsuperscript{(5)} आपन \VUgras{कुकुर} \textsuperscript{(4)}
के साथे निकल पड़ेला।

%%K t08-c2-02-1 p
२. --- ऊ आपन मोट कपड़ा के मजबूत \VUgras{सिकार के पोसाक}
\textsuperscript{(9)} पहिनले बा, गोड़ में चमड़ा के \VUgras{गेटर} बान्हले
बा, जवना से ऊ ओह गीली घास में चल पाई जहाँ \VUgras{बेंग}
\textsuperscript{(1)} छिपल रहेला ; ऊ आपन \VUgras{बंदूक} \textsuperscript{(6)}
आ आपन \VUgras{सिकार-थैला} \textsuperscript{(10)} उठवले बा, आ लऽ, ऊ चल
पड़ल।

%%K t08-c2-03-1 p
३. --- पीछा के जोस ओकरा के एगो अंगूर के खेत के बीच ले आवत बा जहाँ
अंगूर के लवनी जोर पर बा। अंगूर के किसान के लइकन, जे रोज सड़क पर बकरी
चरावेला, आज ओकनी के मनमाना चरे देत बा, आ बुढ़उ \VUgras{बोकरा}
\textsuperscript{(3)} के खूँटा से बान्ह के — जे छूटल रहित त आपन आजादी के
फायदा उठा के जरूर भाग जाइत — ऊ लोग खुद भी एह मजा में हिस्सा बँटावत बा।
एगो \VUgras{अंगूर के टोकरी} \textsuperscript{(12)} से अंगूर के गुच्छा उठा
लेत बा आ \VUgras{रस} \textsuperscript{(14)} \VUgras{कटोरी}
\textsuperscript{(13)} में टपका के मजा लेत बा, ताकि बिना सड़ल अंगूर के रस
बने। दोसरका आपन माथ पर ऊ \VUgras{पत्ता के मउर} \textsuperscript{(15)}
पहिन लेत बा जवन ऊ अभिए गूँथले बा।

ओह लोग के लगे, बेसरम \VUgras{गौरइया} \textsuperscript{(16)} कोल्हू तक आ
जात बाड़ी सँ अंगूर के दाना चुगे।

%%K t08-c2-04-1 p
४. --- जब लइकन खेलत बा, तब मरद लोग काम करत बा। \VUgras{अंगूर तोड़े
वाला} \textsuperscript{(18)} \VUgras{पीठ के टोकरी} \textsuperscript{(17)} में
अंगूर तहखाना में ले जात बा ; \VUgras{पीपा बनावे वाला} \textsuperscript{(19)}
\VUgras{पीपा} \textsuperscript{(20)} ठीक करत बा, देखत बा कि \VUgras{पीपा
के फट्टी} \textsuperscript{(22)} आ \VUgras{पेंदी} \textsuperscript{(21)}
दुरुस्त बा कि ना, आ टूटल \VUgras{लोहा के घेरा} \textsuperscript{(23)} बदलत
बा। ऊ ओह \VUgras{नाद} \textsuperscript{(25)} पर भी काम कइले बा जवना में
\VUgras{अंगूर} \textsuperscript{(26)} डाल के सड़ावल जाला।

%%K t08-c2-05-1 p
५. --- जब सड़ल पूरा हो जाला, त सराब छान लेहल जाला आ बचल भाग
\VUgras{कोल्हू} \textsuperscript{(28)} पर धइल जाला, आ धीरे-धीरे
\VUgras{पेंच} \textsuperscript{(29)} से दबा के रस निचोड़ल जाला, जवन
\VUgras{कठौत} \textsuperscript{(32)} में बहेला आ ओहसे फेरु \VUgras{सराब}
\textsuperscript{(31)} मिलेला।

%%K t08-tit-8 sub
\VUsaut{6.0mm}
\VUpk{10.2pt}{40}{बीयर आ साइडर।}
\VUsaut{3.4mm}

%%K t08-c2-06-1 p
६. --- \VUgras{तहखाना} \textsuperscript{(27)} के भीत पर \VUgras{हॉप के
लता} \textsuperscript{(30)} के डाढ़ चढ़ल बा। इहाँ ऊ खाली सजावट के पौधा ह,
बाकिर कुछ देस में, जहाँ अंगूर बहुते कम होला, हॉप के खेती \VUgras{बीयर}
बनावे खातिर होला।

%%K t08-c2-07-1 p
७. --- छोट \VUgras{सेब के गाछ} \textsuperscript{(33)} सेब से लदल बा,
जवना से पाकला पर बढ़िया \VUgras{साइडर} बनी। आपन \VUgras{पिंजरा}
\textsuperscript{(34)} में \VUgras{कनारी} \textsuperscript{(35)} आ
\VUgras{तिलियर} \textsuperscript{(36)} जलन भरल आँख से ओह गौरइयन के देखत
बाड़ें जे खुल के उछल-कूद करत बाड़ी सँ।

%%K t08-c2-08-1 p
८. --- जहाँ तक नजर जाला, पहाड़ी के ढलान पर \VUgras{अंगूर के खूँटा}
\textsuperscript{(40)} के कतार लागल बा जवन ऊ सुंदर \VUgras{अंगूर के ठूँठ}
\textsuperscript{(39)} के थामले बा, जवन भारी \VUgras{गुच्छा} से लदल बा।

%%K t08-c2-08-2 p
पूरा साल \VUgras{अंगूर के किसान} \textsuperscript{(42)} आपन \VUgras{अंगूर
के बगइचा} \textsuperscript{(38)} के जतन में लागल रहलें आ अबहीं ओकर मेहनत के
इनाम सुंदर फसल से मिलत बा। \VUgras{अंगूर तोड़े वाली}
\textsuperscript{(41)} ओह नाद के भरत बाड़ी जवन \VUgras{खच्चर}
\textsuperscript{(43)} से खींचल गाड़ी पर धइल बा, आ सब लोग खुस बा।

%%K t08-tit-9 sub
\VUsaut{5.0mm}
\VUpk{10.2pt}{40}{मछरी मारल।}
\VUsaut{3.4mm}

%%K t08-c2-09-1 p
९. --- ओहिजे हाल \VUgras{बंसी वाला} \textsuperscript{(44)} लोग के बा, जे
भोरही से आपन \VUgras{बंसी} \textsuperscript{(45)} लेके पानी के किनारे आ
जमल बाड़ें।

%%K t08-c2-09-2 p
\VUgras{मछरी} \textsuperscript{(47)} \VUgras{काँटा} पर झपटेली जइसहीं ऊ
लोग आपन \VUgras{डोर} \textsuperscript{(46)} पानी में डालेला, आ ओह लोग के
\VUgras{चारा} बदले के फुरसत मुसकिल से मिलेला कि नया सिकार डोर के छोर पर
फड़फड़ाए लागेला।

%%K t08-c2-09-3 p
तबो \VUgras{जाल वाला} \textsuperscript{(48)}, जे आपन \VUgras{डोंगी}
\textsuperscript{(49)} से \VUgras{नदी} \textsuperscript{(50)} के बीच में
\VUgras{जाल} फेंकेला, जादे मछरी पकड़ेला।

%%K t08-c2-10-1 p
१०. --- पतझड़ नीक घूमे के मौसम ह : पैदल, \VUgras{घोड़ा पर}
\textsuperscript{(52)}, \VUgras{साइकिल से} \textsuperscript{(53)},
\VUgras{मोटर से} \textsuperscript{(54)}। \VUgras{सड़क} \textsuperscript{(51)}
साफ बा आ लोग आखिरी नीक दिन के फायदा उठावेला, काहेकि लऽ, \VUgras{अबाबील}
\textsuperscript{(56)} पहिलही से छत पर जुटे लागल बाड़ी सँ ताकि पूरा पाँख
फइला के जादे गरम देस उड़ जासु। पुरनका \VUgras{अखरोट के गाछ}
\textsuperscript{(57)} के पत्ता पियर पड़त बा, \VUgras{अखरोट} गिरत बा, आ
ऊँच \VUgras{गाछ} जे \VUgras{गुच्छा} \textsuperscript{(58)} नियर
\VUgras{टीला} \textsuperscript{(59)} पर जुटल बा जल्दिए पतझड़ में नंगा हो जाई।

%%K t08-c2-11-1 p
११. --- \VUgras{सूरज} \textsuperscript{(60)} देर से उगे लागल बा, दिन छोट
होत बा, \VUgras{भोर} के हल्का चुभत ठंढ महसूस होखे लागल बा, आ लऽ,
\VUgras{चाहा} \textsuperscript{(61)} के झुंड पहिलही से हमनी के एह मेहमान-बिरोधी
मौसम के छोड़त बाड़ें।
